\par\noindent\SourceLeaf{53}%
und dem allgemeinen Glied $h^{\text{tus}}$

\[
= \frac{1}{2^{2q+2h-2}} {}^{2q+2h-1}\mathfrak{B}^{(2q+h-1)}  .  {}^{\nu}\mathfrak{B}^{(2q+2h-1)} n^{2q+2h-1}.
\]

C) Auf dieselbe Weise findet man das Gesetz für den Koeffizienten zum Kosinus eines geraden Vielfachen, etwa zu $\cos 2q\varphi$. Denn der Kosinus eines solchen Vielfachen kommt nur in geraden Potenzen von $\cos\varphi$ vor, angefangen bei derjenigen mit dem Exponenten $2q$. Das Anfangsglied der Reihe für den Koeffizienten von $\cos 2q\varphi$ wird daher aus dem Glied der Grundreihe genommen, das $\cos\varphi^{2q}$ enthält, also aus ${}^{\nu}\mathfrak{B}^{(2q)} n^{2q}\cos\varphi^{2q}$, und zwar aus dessen $q^{\text{to}}$ Glied. Es ergibt sich also (da für $\cos\varphi^{2q}$ das $q^{\text{tus}}$ Glied lautet

\[
= \frac{1}{2^{2q-1}} {}^{2q}\mathfrak{B}^{(2q)} \cos 2q\varphi = \frac{1}{2^{2q-1}} \cos 2q\varphi) \qquad = \frac{1}{2^{2q-1}} {}^{\nu}\mathfrak{B}^{(2q)} n^{2q}.
\]

\noindent Allgemein ergibt sich aus dem $h^{\text{to}}$ geraden Glied der Grundreihe nach diesem ersten, also aus der Größe ${}^{\nu}\mathfrak{B}^{(2q+2h)} n^{2q+2h} \cos\varphi^{2q+2h}$, der zu $\cos 2q\varphi$ gehörende $h^{\text{ta}}$ Teil des Koeffizienten. Entwickelt man nun $\cos\varphi^{2q+2h}$, so ist dessen $q^{\text{tus}}$ Glied

\[
= \frac{1}{2^{2q+2h-1}} {}^{2q+2h}\mathfrak{B}^{(2q+h)} \cos 2q\varphi.
\]

\noindent So ergibt sich für den zu $\cos 2q\varphi$ gehörenden Koeffizienten als $h^{\text{tus}}$ Glied nach dem ersten

\[
= \frac{1}{2^{2q+2h-1}} {}^{2q+2h}\mathfrak{B}^{(2q+h)}  .  {}^{\nu}\mathfrak{B}^{(2q+2h)} n^{2q+2h}.
\]

D) Die gestellte Aufgabe ist damit tatsächlich gelöst. Für die aus $(1+n\cos\varphi)^{\nu}$ gewonnene Reihe, die nach der Ersetzung jeder Potenz durch eine Reihe von Kosinus von Vielfachen entsprechend geordnet wurde, ist das Anfangsglied ordnungsgemäß angegeben (vgl. A). Ebenso sind für beide Arten von Gliedern, die in der Reihe vorkommen können, das Anfangsglied und das allgemeine Glied des jeweiligen Koeffizienten ordnungsgemäß angegeben. Die einen enthalten den Kosinus eines ungeraden Vielfachen, $\cos(2q-1)\varphi$ (vgl. B), die anderen den Kosinus eines geraden Vielfachen, $\cos 2q\varphi$ (vgl. C). Diese Darstellung des gefundenen Gesetzes lässt sich aber durch die Beobachtung kürzen, dass für die Koeffizienten der Kosinus gerader Vielfacher genau dasselbe Gesetz gilt wie für die der Kosinus ungerader Vielfacher. Das sieht man wie folgt. Das Anfangsglied des zu $\cos(2q-1)\varphi$ gehörenden Koeffizienten war

\[
\frac{1}{2^{2q-2}} {}^{\nu}\mathfrak{B}^{(2q-1)} n^{2q-1},
\]

\noindent und das allgemeine Glied $h^{\text{tus}}$ war

\[
\frac{1}{2^{2q+2h-2}} {}^{2q+2h-1}\mathfrak{B}^{(2q+h-1)}  .  {}^{\nu}\mathfrak{B}^{(2q+2h-1)} n^{2q+2h-1}.
\]
